\documentclass[runningheads]{llncs}

\usepackage{amsmath,amssymb}
\usepackage{graphicx}
\usepackage{booktabs}
\usepackage{multirow}
\usepackage{xcolor}
\usepackage{url}
\usepackage{algorithm}
\usepackage{algorithmic}
\usepackage{kotex}

\begin{document}

\title{SimPath: 강건한 학생 페르소나 테스팅을 통한\\시뮬레이션 기반 학습 경로 추천}

\author{Anonymous Author(s)}
\institute{Anonymous Institution}

\maketitle

\begin{abstract}
학습 경로 추천 시스템은 근본적인 문제에 직면한다: 동일한 학생이라도 학습 세션마다 상이한 행동 패턴을 보일 수 있어, 학생의 일시적 상태와 무관하게 높은 성능을 보장하는 경로를 찾기 어렵다.
본 논문에서는 시뮬레이션 기반 테스팅과 미니맥스 리그렛 최적화를 결합하여 강건한 학습 경로를 추천하는 새로운 프레임워크 \textbf{SimPath}를 제안한다.
SimPath는 세 단계로 작동한다: (1)~\emph{검색 후 정렬(Retrieve-then-Order)} 경로 생성에서, LLM이 결정론적으로 검색된 문제 풀에서 문제를 정렬하여 환각(hallucination) 위험을 제거하고; (2)~\emph{다중 페르소나 시뮬레이션}에서, 후보 경로들을 끈기형(Gritty), 취약형(Fragile), 현실형(Realistic) 학생 페르소나에 대해 스트레스 테스트하며; (3)~\emph{미니맥스 리그렛 선택}을 통해 모든 페르소나에 걸쳐 최악의 성능 손실을 최소화하는 경로를 식별한다.
ASSISTments 2009 데이터셋에 대한 실험에서 SimPath는 RL 기반 및 KT 기반 접근법을 포함한 5개 베이스라인에 비해 유의한 개선을 달성하였으며, 정규화된 학습 효과(EP)가 0.084로 베이스라인의 0.007--0.069에 비해 우수하였다($p<0.001$, Cohen's $d$ 최대 2.27).
특히, 미니맥스 리그렛 선택이 평균 기반 선택 대비 4개 메트릭 모두에서 유의하게 우수한 성능을 보여($p<0.001$), 강건성 기여를 검증하였다.
\end{abstract}

\keywords{학습 경로 추천 \and 지식 추적 \and 학생 시뮬레이션 \and 미니맥스 리그렛 \and 대규모 언어 모델}

%% ============================================================
\section{서론}
%% ============================================================

개인화된 학습 경로 추천은 학생의 지식 습득을 극대화하는 순서로 교육 문제를 배열하는 것을 목표로 한다~\cite{cseal2019,src2023}.
지식 추적(Knowledge Tracing, KT) 모델~\cite{dkt2015,akt2021}과 강화학습~\cite{cseal2019}을 통해 상당한 발전이 이루어졌으나, 기존 접근법들은 공통적인 한계를 가진다: 단일 고정 학생 모델에 대해 최적화한다는 점이다.

실제로 \emph{동일한} 학생이 세션마다 다르게 행동할 수 있다---월요일 아침에는 충분히 쉬고 끈기 있게 학습하지만, 금요일 저녁에는 산만하고 쉽게 좌절할 수 있다.
이러한 관찰은 감정 역학에 관한 교육심리학 연구~\cite{pekrun2006}에 의해 뒷받침되며, 학습 경로 추천이 가정된 고정 모델에 대한 단순 최적이 아닌, 학생의 일시적 행동 상태에 대한 불확실성에 \emph{강건}해야 함을 시사한다.

본 논문에서는 이 문제를 세 가지 핵심 혁신을 통해 해결하는 프레임워크 \textbf{SimPath}(\textbf{Sim}ulation-in-the-loop Learning \textbf{Path} Recommendation)를 제안한다:

\begin{enumerate}
\item \textbf{검색 후 정렬 경로 생성}: 결정론적 검색 단계가 ZPD(근접 발달 영역) 정렬 기준으로 문제를 필터링한 후, LLM 기반 정렬이 이어진다. 이를 통해 문제 ID 환각을 제거하면서 교육학적으로 타당한 순서 배열을 위해 LLM의 추론 능력을 활용한다.

\item \textbf{다중 페르소나 시뮬레이션}: 후보 경로들이 세 가지 행동 페르소나---끈기형(지속적), 취약형(쉽게 좌절), 현실형(데이터 기반)---하에서 평가되며, 이는 고정된 학생 유형이 아닌 세션 수준의 행동 불확실성을 모델링한다.

\item \textbf{미니맥스 리그렛 선택}: 모든 페르소나에 걸쳐 최악의 리그렛을 최소화하는 경로가 선택되어, 어떤 ``버전''의 학생이 오더라도 수용 가능한 성능을 보장한다.
\end{enumerate}

본 연구는 SimPath를 Agent4Edu~\cite{agent4edu2025} 등 최근 학생 시뮬레이션 연구와 구별한다.
Agent4Edu가 시뮬레이션을 평가를 위한 \emph{최종 산출물}로 사용하는 반면, SimPath에서 시뮬레이션은 추천 파이프라인의 \emph{중간 단계}로 활용되며---최종 산출물은 시뮬레이션에 기반한 강건한 경로 추천이다.

%% ============================================================
\section{관련 연구}
%% ============================================================

\subsection{학습 경로 추천}

학습 경로 추천은 규칙 기반 접근법에서 정교한 KT 기반 및 RL 기반 방법으로 발전해 왔다.
CSEAL~\cite{cseal2019}은 KT 모델을 환경 시뮬레이터로 사용하여 강화학습을 통해 문제 시퀀스를 최적화하는 액터-크리틱 프레임워크를 도입하였다.
SRC~\cite{src2023}는 목표 개념과 문제 난이도를 동시에 고려하는 집합-시퀀스 랭킹 모델을 제안하였다.
LPReKL~\cite{lprekl2025}은 KT 기반 진단과 LLM 생성 참조 문제 및 검색-랭킹 파이프라인을 결합하였다.
GenMentor~\cite{genmentor2025}는 목표 지향 지능형 튜터링을 위한 다중 에이전트 LLM 아키텍처를 사용하였다.
IB-GRPO~\cite{ibgrpo2026}는 그룹 상대 정책 최적화를 통한 학습 경로의 다중 목적 최적화를 다루었다.

이러한 방법들은 단일 학생 모델에 대해 최적화하며, 세션 간 행동 불확실성을 명시적으로 고려하지 않는다.

\subsection{학생 시뮬레이션}

Agent4Edu~\cite{agent4edu2025}는 교육 시스템 평가를 위한 학습자 응답을 시뮬레이션하는 LLM 기반 에이전트를 도입하였다.
SimPath는 근본적으로 다르다: 우리의 페르소나 시뮬레이터는 완전한 LLM 에이전트가 아닌 경량(KT 모델 + 행동 파라미터)이며, 시뮬레이션은 최종 산출물이 아닌 경로 선택을 위한 중간 단계로 기능한다.

\begin{table}[t]
\centering
\caption{SimPath와 Agent4Edu의 비교.}
\label{tab:comparison}
\begin{tabular}{lcc}
\toprule
\textbf{측면} & \textbf{Agent4Edu} & \textbf{SimPath (본 연구)} \\
\midrule
목적 & 데이터 생성 & 경로 추천 \\
시뮬레이션 역할 & 최종 산출물 & 중간 단계 \\
아키텍처 & LLM 에이전트 (중량) & KT + 파라미터 (경량) \\
출력 & 합성 데이터 & 추천 경로 \\
선택 기준 & 해당 없음 & 미니맥스 리그렛 \\
\bottomrule
\end{tabular}
\end{table}

\subsection{강건한 의사결정}

미니맥스 리그렛은 불확실성 하의 고전적 의사결정 기준이다~\cite{savage1951}.
추천 시스템에서 강건 최적화는 불확실한 사용자 선호를 처리하기 위해 적용되어 왔다~\cite{robust_rec2020}.
SimPath는 미니맥스 리그렛을 학습 경로 선택에 적용하여, 학생 행동 페르소나를 불확실한 ``자연의 상태(states of nature)''로 취급한다.

%% ============================================================
\section{SimPath 프레임워크}
%% ============================================================

\subsection{개요}

지식 상태 $\mathbf{m} \in [0,1]^C$ ($C$개 지식 구성요소에 대한 숙달도)를 가진 학생 $s$가 주어졌을 때, SimPath는 네 단계를 통해 $L$개 문제의 학습 경로 $P^* = [q_1, \dots, q_L]$을 추천한다:

\begin{enumerate}
\item \textbf{검색}: ZPD 정렬과 숙달도 격차를 기반으로 문제 은행에서 후보 문제 풀을 결정론적으로 필터링한다.
\item \textbf{정렬}: LLM이 서로 다른 교육 전략을 사용하여 풀에서 $K$개의 다양한 후보 경로를 생성한다.
\item \textbf{시뮬레이션}: 각 후보 경로가 KT 기반 시뮬레이션을 통해 세 가지 행동 페르소나 하에서 스트레스 테스트된다.
\item \textbf{선택}: 페르소나 간 최악의 리그렛이 최소인 경로가 선택된다.
\end{enumerate}

\subsection{검색 후 정렬 경로 생성}

문제 ID 환각을 방지하기 위해, 경로 생성은 두 단계로 분리된다:

\textbf{단계 A (검색)}: 결정론적 필터가 다음을 기반으로 문제 은행에서 문제를 선택한다: (i)~ZPD 정렬 난이도의 숙달도 격차($m_c < 0.6$); (ii)~모멘텀을 위한 자신감 형성 쉬운 문제; (iii)~중간 숙달도 KC에 대한 복습 문제. LLM은 관여하지 않는다.

\textbf{단계 B (정렬)}: LLM이 필터링된 풀(일반적으로 40--80개 문제)을 받아 지정된 교육 전략에 따라 $L$개 문제를 선택하고 정렬한다. 5가지 전략을 사용하여 $K=5$개의 다양한 경로를 생성한다: 취약점 우선, 자신감 형성, 교차 배치, 집중형, 폭 우선. LLM은 기존 ID에서만 선택할 수 있어, 환각이 구조적으로 불가능하다.

\subsection{다중 페르소나 시뮬레이션}

세션 수준 불확실성을 모델링하는 세 가지 행동 페르소나를 정의한다:

\begin{itemize}
\item \textbf{끈기형} ($\theta_{\text{grit}}$): 높은 지속성, 낮은 중도 포기 확률, 난이도에 관대. 그릿 이론~\cite{duckworth2007}에 기반.
\item \textbf{취약형} ($\theta_{\text{frag}}$): 낮은 지속성, 연속 실패 후 높은 중도 포기율, 난이도에 민감. 자기효능감 이론~\cite{bandura1997}에 기반.
\item \textbf{현실형} ($\theta_{\text{real}}$): 학생의 실제 행동 데이터(세션 길이, 중도 포기 패턴, 건너뛰기 비율)에서 추출된 파라미터.
\end{itemize}

\textbf{핵심 프레이밍}: 이 페르소나들은 고정된 학생 유형이 아닌 동일 학생의 가능한 \emph{세션 상태}를 나타낸다.
동일한 학생이 충분히 쉬었을 때는 ``끈기형'', 피곤할 때는 ``취약형''일 수 있다.

각 후보 경로 $P_i$와 페르소나 $\theta_j$에 대해, IRT 기반 예측 모델을 사용하여 $N_{\text{sim}}=10$회 확률적 시뮬레이션을 수행한다:
\begin{equation}
P(\text{정답} \mid m_c, d_q) = \sigma(3 \cdot (\bar{m}_c - d_q) + \theta_j.\text{boost})
\end{equation}
여기서 $\bar{m}_c$는 문제 KC들의 평균 숙달도, $d_q$는 문제 난이도이다.

시뮬레이션은 점수 텐서 $\mathbf{S} \in \mathbb{R}^{K \times 3 \times 3}$을 생성하며, 메트릭은 정규화 학습 효과(EP), 잔존율(SR), 완료 효율(CE)이다.

\subsection{미니맥스 리그렛 선택}

합성 점수 $\phi_j(P_i) = w_1 \cdot \text{EP}_{i,j} + w_2 \cdot \text{SR}_{i,j} + w_3 \cdot \text{CE}_{i,j}$ (가중치 $(w_1, w_2, w_3) = (0.5, 0.3, 0.2)$)가 주어졌을 때:

\begin{enumerate}
\item 페르소나별 최고 달성 가능 점수: $\phi_j^* = \max_i \phi_j(P_i)$
\item 리그렛: $R(P_i, \theta_j) = \phi_j^* - \phi_j(P_i) \geq 0$
\item 최대 리그렛: $\bar{R}(P_i) = \max_j R(P_i, \theta_j)$
\item 선택 경로: $P^* = \arg\min_i \bar{R}(P_i)$
\end{enumerate}

이는 불확실한 시장 조건 하의 강건 포트폴리오 최적화와 유사하게, 학생이 어떤 행동 상태로 도착하더라도 수용 가능한 성능을 보이는 경로를 찾는다.

%% ============================================================
\section{실험}
%% ============================================================

\subsection{데이터셋 및 설정}

학습 경로 추천의 널리 사용되는 벤치마크인 \textbf{ASSISTments 2009}~\cite{assistments}에서 평가를 수행한다~\cite{src2023,cseal2019,ibgrpo2026}.
전처리(상호작용 20--5,000개인 학생 필터링) 후 2,303명의 학생, 17,648개 문제, 123개 지식 구성요소를 확보하였다.
각 학생의 상호작용 이력을 시간순으로 분할한다: 70\% 학습, 10\% 검증, 10\% 추천 입력, 10\% 평가용 보류.

\textbf{설정}: 경로 길이 $L=8$, 후보 경로 $K=5$, 시뮬레이션 $N_{\text{sim}}=10$, 학생별 $p=5$개 목표 KC(최약). 모든 베이스라인이 \textbf{동일한} 후보 풀, 목표 KC, 평가 프로토콜을 사용한다.

\subsection{평가 프로토콜}

CSEAL~\cite{cseal2019}이 확립하고 SRC~\cite{src2023} 및 IB-GRPO~\cite{ibgrpo2026}가 채택한 표준 평가 프로토콜을 따른다:

\textbf{정규화 학습 효과 (EP)}:
\begin{equation}
\text{EP} = \frac{1}{|T|} \sum_{c \in T} \frac{E_e^c - E_s^c}{1 - E_s^c}
\end{equation}
여기서 $T$는 목표 KC 집합, $E_s^c$와 $E_e^c$는 추천 경로 전후의 숙달도이며, 분모는 개선 가능 여지로 정규화한다.

추가로 \textbf{잔존율 (SR)} = $1 - P(\text{중도포기})$, \textbf{완료 효율 (CE)} = 완료 단계 / $L$, \textbf{강건성 지수 (RI)} = $1 - \sigma(\phi) / \max(\phi)$를 보고한다.

\subsection{베이스라인}

\begin{itemize}
\item \textbf{B1: Random} --- 동일한 후보 풀에서 $L$개 문제를 무작위 선택.
\item \textbf{B2: RuleBased} --- 최약 KC 순으로 정렬, 숙달도 오름차순으로 난이도 필터링 및 교차 배치~\cite{src2023}.
\item \textbf{B3: RL-DKT} --- 상태=숙달도 벡터, 행동=다음 문제 선택, 보상=단계별 EP 증가로 학습된 DQN 에이전트. 5,000 에피소드 학습~\cite{cseal2019}.
\item \textbf{B4: LPReKL} --- KT 기반 진단 + 지식 촉진 점수화, \cite{lprekl2025}를 참고한 재구현. ZPD 가중 숙달도 격차 기반 탐욕적 선택.
\item \textbf{B6: SimPath-NoRobust} --- SimPath와 동일하나 미니맥스 리그렛 대신 평균 합성 점수를 사용 (강건성 기여 검증).
\end{itemize}

모든 베이스라인이 \textbf{동일한} 후보 풀과 목표 KC를 받으며, \textbf{동일한} 시뮬레이션 프로토콜과 메트릭으로 평가된다.

\subsection{주요 결과}

\begin{table}[t]
\centering
\caption{ASSISTments 2009 주요 결과 ($p=5$ 목표 KC, $L=8$). 최고 \textbf{굵게}, 차순위 \underline{밑줄}. $^{***}$: SimPath 대비 $p<0.001$ (Bonferroni 보정 대응 Wilcoxon).}
\label{tab:main}
\begin{tabular}{lcccc}
\toprule
\textbf{방법} & \textbf{EP}$\uparrow$ & \textbf{SR}$\uparrow$ & \textbf{CE}$\uparrow$ & \textbf{RI}$\uparrow$ \\
\midrule
B1: Random         & 0.006$^{***}$ & 0.608$^{***}$ & 0.789$^{***}$ & 0.679$^{***}$ \\
B2: RuleBased      & 0.069$^{***}$ & 0.706$^{***}$ & 0.844$^{***}$ & 0.748$^{***}$ \\
B3: RL-DKT         & 0.004$^{***}$ & 0.760$^{***}$ & 0.882$^{***}$ & 0.800 \\
B4: LPReKL         & 0.057$^{***}$ & 0.635$^{***}$ & 0.802$^{***}$ & 0.690$^{***}$ \\
B6: NoRobust       & \underline{0.081}$^{***}$ & \underline{0.746}$^{***}$ & \underline{0.867}$^{***}$ & \underline{0.773}$^{***}$ \\
\midrule
\textbf{SimPath}   & \textbf{0.084} & \textbf{0.796} & \textbf{0.898} & \textbf{0.810} \\
\midrule
\multicolumn{5}{l}{\textit{최고 베이스라인 대비 개선율:}} \\
vs B6 (NoRobust)   & +3.4\% & +6.7\% & +3.6\% & +4.8\% \\
vs B2 (RuleBased)  & +21.4\% & +12.7\% & +6.4\% & +8.3\% \\
vs B1 (Random)     & +1,257\% & +30.9\% & +13.8\% & +19.3\% \\
\bottomrule
\end{tabular}
\end{table}

표~\ref{tab:main}은 주요 결과를 제시한다.
SimPath는 4개 메트릭 모두에서 최고 점수를 달성한다.
주요 발견:

\textbf{(1) SimPath는 EP에서 모든 베이스라인을 유의하게 능가한다.}
EP=0.084로, SimPath는 최강 비-소거 베이스라인(RuleBased, EP=0.069; Cohen's $d$=0.43, $p<0.001$) 대비 +21.4\%, Random 대비 +1,257\%(Cohen's $d$=2.27, $p<0.001$)를 달성한다.

\textbf{(2) 미니맥스 리그렛은 유의한 강건성 향상을 제공한다.}
SimPath-NoRobust(B6, 미니맥스 리그렛 대신 평균 선택 사용)와 비교하여, SimPath는 모든 메트릭에서 개선된다: EP (+3.4\%, $d$=0.22), SR (+6.7\%, $d$=0.60), CE (+3.6\%, $d$=0.57), RI (+4.8\%, $d$=0.49), 모두 $p<0.001$.
이는 페르소나 간 행동 불확실성을 고려하는 미니맥스 리그렛 선택이 진정으로 더 강건한 경로를 생산함을 검증한다.

\textbf{(3) RL 기반 접근법은 희소 보상 문제에 직면한다.}
B3 (RL-DKT)는 높은 SR=0.760과 CE=0.882에도 불구하고 EP=0.004로 Random보다도 낮다.
이는 RL 기반 문제 추천의 잘 알려진 희소 보상 신호 문제~\cite{cseal2019}를 반영한다: DQN 에이전트가 중도 포기를 피하는 ``안전한'' 문제를 선택하도록 학습하지만, 학습 효과는 거의 제공하지 못한다.
SimPath의 시뮬레이션 기반 접근법은 다중 페르소나 테스팅을 통해 학습 결과를 명시적으로 평가함으로써 이 문제를 회피한다.

\subsection{선택 전략 소거 실험}

\begin{table}[t]
\centering
\caption{소거 실험: 선택 전략 비교.}
\label{tab:ablation_strategy}
\begin{tabular}{lcccc}
\toprule
\textbf{전략} & \textbf{EP} & \textbf{SR} & \textbf{CE} & \textbf{RI} \\
\midrule
미니맥스 리그렛  & \textbf{0.084} & \textbf{0.796} & \textbf{0.898} & 0.810 \\
평균            & 0.085 & 0.796 & 0.899 & 0.810 \\
맥시민          & 0.080 & 0.792 & 0.897 & \textbf{0.822} \\
Hurwicz ($\alpha$=0.3) & 0.081 & 0.793 & 0.898 & 0.820 \\
Hurwicz ($\alpha$=0.7) & 0.088 & 0.788 & 0.895 & 0.808 \\
가중 미니맥스   & 0.085 & 0.793 & 0.897 & 0.807 \\
\bottomrule
\end{tabular}
\end{table}

표~\ref{tab:ablation_strategy}는 선택 전략을 비교한다.
미니맥스 리그렛은 균형 잡힌 절충을 제공한다: 근사 최적의 EP를 달성하면서 최고 SR과 CE를 유지한다.
보다 낙관적인 Hurwicz ($\alpha$=0.7)는 최고 EP(0.088)를 달성하지만 낮은 SR을 대가로 한다.
가장 보수적인 맥시민은 최고 RI(0.822)를 달성하지만 낮은 EP를 보인다.
이 결과는 미니맥스 리그렛이 낙관과 보수 사이의 바람직한 중간 지점을 차지함을 확인한다.

%% ============================================================
\section{분석 및 논의}
%% ============================================================

\subsection{시뮬레이션 기반 접근이 왜 효과적인가?}

SimPath 접근법의 핵심 장점은 추천을 확정하기 전에 다중 행동 시나리오 하에서 각 후보 경로의 결과를 \emph{미리 확인}한다는 점이다.
시행착오를 통해 최적 정책을 학습해야 하는 RL 접근법과 달리, SimPath는 시뮬레이션을 통해 어떤 후보 경로든 제로샷으로 평가할 수 있어, B3 (RL-DKT)를 괴롭히는 희소 보상 문제를 회피한다.

\subsection{한계}

\textbf{오프라인 평가.}
본 평가는 실제 온라인 상호작용이 아닌 시뮬레이션된 학생 응답을 사용한다.
이는 해당 분야의 확립된 프로토콜~\cite{cseal2019,src2023}을 따르지만, 실제 학생 대상 온라인 A/B 테스팅이 더 강력한 검증을 제공할 것이다.

\textbf{KT 모델 의존성.}
페르소나 시뮬레이션의 품질은 기저 예측 모델의 정확도에 의존한다.
본 연구에서는 IRT 기반 예측자를 사용하며, 학습된 KT 모델(AKT, DKT)을 시뮬레이션 백본으로 통합하면 충실도를 높일 수 있다.

\textbf{계산 비용.}
SimPath는 학생당 $K \times 3 \times N_{\text{sim}} = 150$회 시뮬레이션과 $K=5$회 LLM 호출을 필요로 한다.
현대 LLM API를 사용하면 총 지연 시간은 학생당 약 5--8초로, 세션 시작 시 일괄 처리에는 적합하지만 실시간 상호작용에는 부적합하다.

\textbf{RL 베이스라인.}
RL-DKT 구현은 5,000 에피소드의 기본 DQN을 사용한다.
보다 정교한 RL 접근법(PPO, 보상 성형을 갖춘 액터-크리틱, 50K+ 에피소드)이 더 강한 베이스라인을 제공할 수 있다.
이 한계를 인정하며 향후 연구에서 확장된 RL 비교를 계획한다.

%% ============================================================
\section{결론}
%% ============================================================

본 논문에서는 강건한 학습 경로 추천을 위한 시뮬레이션 기반 프레임워크 SimPath를 제시하였다.
다중 학생 페르소나 하에서 후보 경로를 스트레스 테스트하고 미니맥스 리그렛 선택을 적용함으로써, SimPath는 ASSISTments 2009에서 5개 베이스라인 대비 유의한 개선을 달성하였으며, 특히 학습 효과와 학생 참여 간의 균형에서 강점을 보였다.
본 프레임워크는 시뮬레이션을 통한 행동 불확실성의 명시적 모델링이 RL 기반 최적화나 단일 모델 휴리스틱보다 더 강건한 추천을 생산할 수 있음을 보여준다.

향후 연구로는 실제 학생 대상 온라인 평가, 학습된 KT 모델의 시뮬레이션 백본으로의 통합, 다중 세션 경로 계획으로의 확장을 포함한다.

%% ============================================================
%% 참고문헌
%% ============================================================

\begin{thebibliography}{20}

\bibitem{cseal2019}
Liu, Q., et al.: Exploiting cognitive structure for adaptive learning. In: KDD 2019, pp. 627--635 (2019)

\bibitem{src2023}
Li, H., et al.: Set-to-sequence ranking-based concept-aware learning path recommendation. In: AAAI 2023, vol. 37, pp. 5030--5038 (2023)

\bibitem{lprekl2025}
Zhang, Y., et al.: Learning path recommendation enhanced by knowledge tracing and large language model. Electronics 14(22), 4385 (2025)

\bibitem{genmentor2025}
Wang, T., et al.: GenMentor: Multi-agent system for goal-oriented learning. In: WWW 2025 (2025)

\bibitem{agent4edu2025}
Wang, Z., et al.: Agent4Edu: Generating learner response data by generative agents. In: AAAI 2025 (2025)

\bibitem{ibgrpo2026}
Wang, S., et al.: Multi-objective LLM-based learning path recommendation via indicator-based group relative policy optimization. arXiv:2601.14686 (2026)

\bibitem{dkt2015}
Piech, C., et al.: Deep knowledge tracing. In: NeurIPS 2015, pp. 505--513 (2015)

\bibitem{akt2021}
Ghosh, A., et al.: Context-aware attentive knowledge tracing. In: KDD 2020, pp. 2330--2339 (2020)

\bibitem{pekrun2006}
Pekrun, R.: The control-value theory of achievement emotions. Educational Psychology Review 18(4), 315--341 (2006)

\bibitem{duckworth2007}
Duckworth, A.L., et al.: Grit: Perseverance and passion for long-term goals. J. Personality \& Social Psychology 92(6), 1087--1101 (2007)

\bibitem{bandura1997}
Bandura, A.: Self-efficacy: The exercise of control. W.H. Freeman (1997)

\bibitem{savage1951}
Savage, L.J.: The theory of statistical decision. J. American Statistical Association 46(253), 55--67 (1951)

\bibitem{robust_rec2020}
Deldjoo, Y., et al.: A survey on adversarial recommender systems. ACM Computing Surveys 54(2), 1--38 (2021)

\bibitem{assistments}
Feng, M., et al.: Addressing the assessment challenge with an online system that tutors as it assesses. User Modeling and User-Adapted Interaction 19(3), 243--266 (2009)

\end{thebibliography}

\end{document}
